\section*{引言}\label{sec:0}
\addcontentsline{toc}{section}{引言}

谱定理与傅里叶变换在数学与物理中无处不在。就其根本而言，它们关乎如何利用阿贝尔对称性（例如空间平移）把函数分解为简单、单色的构件。人们自然会想：非阿贝尔对称性的对应物想必更为强大、也更为有用。然而，我们今天对非阿贝尔对偶的理解，仍不免让人想起“盲人摸象”的寓言。数论学家已经为这种非阿贝尔对偶发现了一幅宏大图景——朗兰兹纲领，其回声遍及表示论，应用众多（其中最著名者即费马大定理的证明）。物理学家则独立地给出了另一幅宏大图景：麦克斯韦（Maxwell）理论中电磁对偶的非阿贝尔形式，即所谓 Montonen--Olive S-对偶，其回声同样遍及数学物理。得益于对几何朗兰兹对偶（Geometric Langlands Correspondence，以下简称 GLC）的研究，我们现在明白：这两个群体实际上是在勘察同一头巨象，而且如今已能交换观察、协调彼此的考察。事实上，我们或许已经辨认出这头巨兽跳动的心脏——由因子化代数（factorization algebra）这一概念所提供的、隐藏的交换性。

这一探索近来迎来了一个重大里程碑：Gaitsgory 与 Raskin 会同 Arinkin、Beraldo、\mbox{Campbell}、Chen、Faergeman、Lin 与 Rozenblyum，证明了非分歧几何朗兰兹猜想。该证明展开为三篇论文 [GR24a, ABC${}^{+}$24a, GR24b] 与一部两卷本专著 [ABC${}^{+}$24b, CCF${}^{+}$24]，其基础是此前的大量技术工作，包括 Lurie 关于高阶代数的基础 [Lur09, Lur]，以及 Gaitsgory 过去二十年间的许多工作。要理解 GLC 的陈述——更不必说其证明——须得穿透一层又一层的坚韧厚皮，并在众多活动部件之间穿行。但我相信，由此获得的洞见与视角足以回报这番努力，并且指向一个远比眼前所见更为宏大的整体。

我想借此机会作一番综述：GLC 究竟是怎样一种陈述，我们为何应当关心它，以及它为何成立。这绝不是一份关于证明的详尽叙述。我的重点放在那些最令我着迷的部分——它的众多联系与视角，以及它帮助揭示的深层结构。我先在 \S\ref{sec:1} 中讨论谱定理与傅里叶变换，以此铺开整个故事；接着在 \S\ref{sec:2} 中给出 GLC 的具体背景（陈述最终落在 \S\ref{sec:2.5}）。\S\ref{sec:3} 与 \S\ref{sec:4} 讨论两个主要结构——因子化与 Hitchin 系统——它们提示 GLC 可能为真；随后在 \S\ref{sec:5} 中概述证明。最后在 \S\ref{sec:6} 中谈一些应用——急于知道我们忙这一切究竟为了什么的读者，不妨先翻到这里——以及接下来可能的方向。

\textit{Caveat lector}（读者须知）：我始终宁可偏向直觉而牺牲精确，宁可引用综述文章而不引原始文献。在引用相关论文与重要进展方面，我远谈不上全面，甚至算不上合格。我也没有讨论 GLC 宇宙的“第一阶段”（Phase I）——从问题在 [Dri83, Lau87] 中的起源，到 [FGV01, Gai04] 中 $\mathrm{GL}_{n}$ 的 Hecke 本征层的构造（该领域早期的综述见 [Fre04, Gai17a]）——而是聚焦于第二阶段（Phase II）的精华，即范畴化的 GLC（categorical GLC），从 [BD] 直到其最近的证明。\zhnote{“第一阶段／第二阶段”（Phase I / Phase II）是作者本人对几何朗兰兹发展史的分期，本译本照原文保留：Phase I 自问题的提出（[Dri83, Lau87]）延伸到 $\mathrm{GL}_{n}$ 的 Hecke 本征层之构造（[FGV01, Gai04]），其综述见 [Fre04, Gai17a]；Phase II 即“范畴化的 GLC”，以 [BD] 为起点，经 Arinkin--Gaitsgory 引入奇异支撑与 $\IndCoh_{\Nilp}$ 的表述（[AG15]），直至 2024 年非分歧猜想的证明（[GR24a, GR24b, ABC${}^{+}$24a, ABC${}^{+}$24b, CCF${}^{+}$24]）。原文所谓 “the greatest hits of Phase II” 即指这一时期的代表性成果。[来源：A 级，Ben-Zvi 原文所述分期；[BD]、[AG15] 与 [Sch26] 互证]}我热忱推荐 [Sch26]，那是对几何朗兰兹纲领现状一份富有洞见、深入详尽的叙述。

\textit{致谢}：感谢“当前进展通报”（Current Events Bulletin）的组织者 Daniel Erman 与 Bianca Viray。\zhnote{\emph{Current Events Bulletin}（“当前进展通报”）是美国数学会在年度联合数学会议（Joint Mathematics Meetings）上组织的综述报告系列，本文的致谢对象即该系列的当届组织者。[来源：B 级，AMS 会议页面 \url{https://www.ams.org/meetings/lectures/current-events-bulletin}]}同样感谢 Alexandre Afgoustidis、Ron Donagi、Nigel Higson、Masoud Kamgarpour、JiWoong Park 与 Sam Raskin 对早期稿本提出的宝贵意见。本研究得到美国国家科学基金会（National Science Foundation）DMS-2302356 号资助，谨致谢意。
